% !PRINTED 425
\weberpage{425}
If, conversely, any solution $x, y$ of (8) is given, then $4n$ is
divisible by the square of every common divisor of $x$ and $y$. We
denote by $\delta$ the greatest common divisor of $x, y$ whose \gesperrt{square} is
at the same time a divisor of $n$, and which is therefore either equal
to the greatest common divisor of $x, y$ or equal to half of it, and we
set
\[
y = \delta y', \quad n = \delta x', % [sic]
\]
and thus obtain a \gesperrt{proper solution} of (7).

If, then, we put for $-m$ all discriminants for which equation (8) has
solutions at all, and understand by $k$ the \gesperrt{number} of these solutions
(with positive $x$), and further by $h(m)$ the class number of primitive
forms of discriminant $-m$, then the formula results
\weq{9}{\Sigma k\, h(m) = 2\sum \partial,}

if $n$ is not a square,
\[
= 2\sum \partial + \sqrt{n} - 1,
\]
if $n$ is a square.

Here it must still be recalled that in the two exceptional cases $m=3$,
$m=4$ only the third part or the half of the number of proper solutions
of (8) was understood by $k$ in § 115.

We will now do justice to this circumstance by understanding by $k$ the
total number of solutions of (8), but setting
\[
h(3) = \tfrac{1}{3}, \quad h(4) = \tfrac{1}{2}
\]
[instead of $h(3)=1$, $h(4)=1$].

It is further to be observed that in forming the sum (9) the solutions
belonging to the value $n'=1$ have not been counted. This comes into
consideration only in the case where $n$ is a square. Equation (7):
\[
4 = y'^2 + m x'^2
\]
then has only the following proper solutions:
\[
\begin{array}{lllll}
m=3: & x'=1, & y'=\pm 1, & k=2, & h=\tfrac{1}{3},\\
m=4: & x'=1, & y'=0, & k=1, & h=\tfrac{1}{2},
\end{array}
\]
and hence, if one wished to remove this exception, there would still
have to be added to (9):
\[
\Sigma k h, \text{ referred to } n'=1, = \tfrac{1}{2}+\tfrac{2}{3} = \tfrac{7}{6}.
\]
